\section{人机对齐与评价}

讲者进一步提出 ``agent judge'' 的思想：不再依赖人工标注整条 trace，而是用经过偏好优化的 agent 作为 judge，对每一次互动的 plan 进行打分，再把这些偏好信号回流到 generation 过程中。
\begin{table}[H]
\centering
\begin{tabular}{lp{6cm}}
\toprule
指标 & 说明 \\
\midrule
Accuracy & 关键字段（目的地、预算、时间）是否正确且一致 \\
Proactivity & 是否在有限轮数内提问足够的问题来补全上下文 \\
Credibility & 自我一致性与 hallucination 频率；judge 会 penalize 自相矛盾 answer \\
Diversity & 多轮对话中的方案数量与覆盖不同策略 \\
\bottomrule
\end{tabular}
\caption{Agent judge 评价指标}
\end{table}

\begin{knowledgebox}{Agent Judge+Human Alignment}
agent judge 的输出会与 human preference 的 proxy（如 plan optimality、constraint compliance）做 soft alignment，并通过 DPO 反向传导，让 generator 更快收敛到用户认可的行为。
\end{knowledgebox}

\subsection{本章小结}

Agent judge 提供了一个高频反馈环：每次对话后的 reward 既用于 judge 本身的 DPO 更新，也用于 calibrate planner，形成 metrics + training 的闭环。

